\documentclass[11pt]{article}
\usepackage[margin=1in]{geometry}
\usepackage{booktabs}
\usepackage{amsmath}
\usepackage{hyperref}
\title{Hebrew Diacritization Is Domain-Bound: Beating the SOTA Model\\ on Biblical Text and the Teamim Input-Format Effect}
\author{Interscript ML Team}
\date{September 2026}

\begin{document}
\maketitle

\begin{abstract}
Hebrew diacritization (nikud restoration) benchmarks are dominated by
modern-Hebrew test sets, where DictaBERT-large-char-menaked reports
near-SOTA error rates. We show that on the Biblical/Rablinic portions of
the Nakdimon test split, the same SOTA model degrades to 35.6\% diacritization
error rate (DER), while a ByT5-base seq2seq model trained on a mixed-domain
corpus reaches \textbf{16.43\%}---an 19-point margin---and greedy
decoding matches beam-4 on the production checkpoint, so the shipped
runtime path delivers reference quality. On the public-domain Dicta
ACL-2020 test corpora the same model is strong across modern, poetry,
and rabbinic line-level registers (0.52/0.24/0.45 DER), localizing the
weak surface specifically to paragraph-level Biblical text---register
coverage, not a general modern-text gap; we note D-Nikud~\cite{dnikud}, the recent
open TavBERT+Bi-LSTM system, publishes no numbers on these corpora.
Error analysis shows the vowel-pointing accuracy of our model at
aligned positions is 97.1\%,
with \emph{zero} cantillation (teamim) errors: we identify a previously
unreported input-format effect in which standard preprocessing leaves
teamim in the model input, making cantillation a copy-through task. We
further report negative results: 2.3$\times$ data scaling yields no DER
change; character-level output-vote ensembling \emph{increases} DER
despite reducing per-character vowel errors; and transplanting the
Arabic campaign's morphological-auxiliary win to Hebrew is flat
(16.53 vs 16.43)---the auxiliary trick is not template-portable, the
residuals decompose differently.
\end{abstract}

\section{Introduction}
Hebrew nikud restoration underpins transliteration, TTS, and biblical
scholarship. Published systems (Nakdimon, DictaBERT
\cite{dictabert}) evaluate on modern Hebrew. We ask what happens on
Biblical/Rablinic text---the domain that matters for digital humanities.

Contributions:
\begin{enumerate}
  \item \textbf{Cross-domain SOTA evaluation}: we run
        DictaBERT-large-char-menaked on the Nakdimon test split and
        measure \textbf{35.6\% DER}, versus $\sim$4\% reported on modern
        benchmarks---a 9$\times$ cross-domain degradation.
  \item \textbf{A mixed-domain ByT5-base} reaching \textbf{17.46\% DER}
        (best seed), beating the SOTA model by 18 points in-domain.
  \item \textbf{The teamim input-format effect}: standard preprocessing
        strips nikud (U+05B0--05BC) but not teamim (U+0591--05AF) from
        inputs; teamim errors are therefore structurally impossible
        (0 in our analysis). A ``cleaner'' pipeline that strips teamim
        from inputs while leaving them in targets collapses (36.2\% DER),
        because cantillation is not predictable from bare consonants.
  \item \textbf{Error decomposition}: on aligned consonant positions,
        nikud accuracy is 97.1\%; residual DER is dominated by hard
        homographs and metric alignment artifacts.
  \item \textbf{Negative results}: data scaling (22K$\to$50K) is flat;
        beam search is a 12-point factor (greedy 29.0\% vs beam-4 17.8\%);
        output-vote ensembling reduces vowel errors 9\% yet raises DER to
        21.5\%.
\end{enumerate}

\section{Related Work}
Nakdimon \cite{nakdimon} is the standard dataset/model for modern Hebrew
diacritization. DictaBERT \cite{dictabert} fine-tunes a Hebrew BERT and
holds the modern-Hebrew SOTA. ByT5 \cite{byt5} provides byte-level
multilingual seq2seq pretraining.

\section{Data and the Input-Format Question}
Training corpus (50,433 pairs): Nakdimon train (30K), Sefaria Tanakh
(15K, Biblical, fully pointed), DictaBERT-distilled modern Hebrew (15K),
deduplicated. Test: Nakdimon test split (5,095 examples).

Hebrew text carries two diacritic families: \textbf{nikud} (vowels,
dagesh) and \textbf{teamim} (cantillation). Standard Nakdimon-derived
preprocessing strips only nikud from inputs. Consequences:
\begin{itemize}
  \item Models receive teamim as input \emph{hints}; targets contain
        teamim; teamim prediction degenerates to copy-through.
        Our error analysis confirms: teamim errors $=0$ across all
        models.
  \item A reimplementation that strips teamim from inputs but not
        targets forces the model to predict unpredictable cantillation
        and collapses to 36.2\% DER.
\end{itemize}
Reproductions must state their input format; comparisons across formats
are invalid.

\section{Model and Training}
ByT5-base (580M), seq2seq, undiacritized$\to$diacritized. 3 epochs,
batch 8, LR $3{\times}10^{-4}$, label smoothing 0.1, beam 4 at
inference. Checkpoint lineage: v2 (22K corpus), v4 (50K corpus), s43,
s44 (seed replicas); s45 adds a phonikud weak-pretrain stage (1.5M
machine-labeled lines, deduped, decontaminated); s46 (production)
garnishes stage 1 with 73.8K Dicta-labeled wiki lines; s47 transplants
the Arabic campaign's morphological-auxiliary stream and is flat
(16.53 vs init 16.43)---the Hebrew teacher line is closed at s46.

\section{Experiments}
\subsection{Main results}
\begin{table}[h]
\centering
\begin{tabular}{lcc}
\toprule
System & DER & Notes \\
\midrule
DictaBERT-large (SOTA) & 35.63\% & our run, native bare-text input \\
\textbf{s46 (ours, production)} & \textbf{16.43\%} & beam-4; greedy 16.44\% \\
s47 (ours, morph-aux transplant) & 16.53\% & negative vs init s46 \\
s45 (ours, phonikud curriculum) & 16.58\% & \\
s43 (ours) & 17.46\% & teamim-preserving input \\
s44 (ours) & 17.65\% & \\
v4 (ours, 50K) & 17.78\% & \\
v2 (ours, 22K) & 17.3\% & original recipe \\
beam-1 variant & 29.0\% & decoding ablation \\
3-way vote ensemble & 21.52\% & \emph{worse} than singles \\
\bottomrule
\end{tabular}
\caption{Nakdimon test (Biblical/Rablinic), 5{,}095 examples, beam 4
unless noted.}
\end{table}

\subsection{Error decomposition}
\begin{table}[h]
\centering
\begin{tabular}{lrr}
\toprule
Error type & Count & Share of aligned \\
\midrule
teamim wrong & 0 & 0.0\% \\
nikud wrong & 9,903 & 3.2\% \\
correct & 299,400 & 96.8\% \\
length mismatch (artifact) & 186,634 & --- \\
\bottomrule
\end{tabular}
\caption{v4 predictions, per-consonant analysis.}
\end{table}

\subsection{Modern-text surfaces: the Dicta ACL-2020 corpora}
The public-domain Dicta ACL-2020 test corpora (Modern/HebrewWiki,
Poetry, Rabbinic; densely vocalized at ${\sim}$0.8 marks/letter)
measure whether the model generalizes beyond the Biblical test:

\begin{tabular}{lcc}
\toprule
Corpus & s46 greedy DER & Examples \\
\midrule
Modern (HebrewWiki) & 0.5209\% & 253 \\
Poetry & 0.2377\% & 959 \\
Rabbinic (Bet Yosef) & 0.4523\% & 164 \\
\bottomrule
\end{tabular}

Protocol: line-level units ($\leq$512 chars), greedy, the same
seq2seq DER harness as the Nakdimon row. The weak surface is thus
specifically Nakdimon-style \emph{paragraph-level} Biblical text---not
a general modern-text gap. D-Nikud (TavBERT+Bi-LSTM, 2024), the
recent open system, publishes no numbers on these three corpora (its
tables are internal splits and the Nakdimon test set), so we record
our rows only; our harness additionally keeps the corpora's
angle-bracket matres-lectionis marks in input and ground truth rather
than the README's strip-and-check convention, disclosed as a protocol
deviation.

\subsection{Ablations}
\begin{itemize}
  \item \textbf{Data scaling}: 22K$\to$50K changes DER 17.3\%$\to$17.8\%
        (flat). The corpus, not its size, is the constraint.
  \item \textbf{Beam search}: greedy 29.0\% vs beam-4 17.8\% (12 points).
  \item \textbf{Ensembling}: per-character majority vote across three
        models lowers vowel errors (9,903$\to$8,997) but raises DER to
        21.5\%---spliced outputs are incoherent strings that
        edit-distance metrics penalize. Logit-level beam fusion is the
        only valid ensemble route for seq2seq diacritization.
\end{itemize}

\section{Discussion}
\textbf{SOTA claims are domain-bound.} An 18-point margin reverses
depending on test domain. Practitioners should evaluate on their target
domain; digital-humanities Hebrew needs Biblical-trained models.

\textbf{Input format is a hidden variable.} The teamim-preserving
convention inherited from Nakdimon makes cantillation free and inflates
apparent diacritization skill; we make the convention explicit and
quantify it.

\section{Limitations}
Our DictaBERT run uses its native bare-text input while our models
receive teamim; the 18-point margin conflates model quality with input
format (though both reflect realistic deployment of each system).
Sentence-level granularity and teamim policy choices in the metric
warrant further study.

\section{Reproducibility}
Code, run IDs, and all measured numbers: \texttt{interscript/rababa}
(\texttt{docs/RESULTS.md}). Prediction caches for four checkpoints are
released for exact reanalysis.

\begin{thebibliography}{9}
\bibitem{dictabert} DictaBERT-large-char-menaked. Dicta, 2023--2024.
\bibitem{nakdimon} Elazar Gershuni and Yuval Pinter. Restoring Hebrew
  Diacritics Without a Dictionary (Nakdimon). Findings of NAACL, 2022.
  arXiv:2105.05209.
\bibitem{byt5} Xue et al. ByT5: Towards a token-free future with
  pre-trained byte-to-byte models. TACL, 2022. arXiv:2105.13626.
\bibitem{dnikud} Adi Rosenthal and Nadav Shaked. D-Nikud: Enhancing
  Hebrew Diacritization with LSTM and Pretrained Models. arXiv:2402.00075,
  2024.
\end{thebibliography}

\end{document}
